% !TEX TS-program = lualatex+se+debug
% !TEX encoding = UTF-8 Unicode

\documentclass[12pt, twoside, a4paper]{extreport}
% -----------------------------------------------------------------------------------
% ------------------------- USTAWIENIA STRONY--------------------------
% -----------------------------------------------------------------------------------
\usepackage{geometry}
\geometry{
	left=10mm,
	right=10mm,
	top=15mm,
	bottom=15mm,
	marginparwidth=0mm,
	marginparsep=0mm
	}
\newcommand{\wstawplik}[1]{\input{../include/#1}}%skraca kod
\wstawplik{headerofficium}
% elementy zmienne co do headerofficium:
\graphicspath{{../img/}} % katalog obrazów
%\svgpath{{../include/}}
\grechangestaffsize{22}%rozmiar wyniku GregorioTEX
\gresetlastline{ragged}% justowanie do lewej
\grechangestyle{initial}{\greIniZallmanCaps\fontsize{36}{20}\selectfont} 
\grechangedim{commentaryraise}{4pt}{scalable} % podnosi komentarz wystarczająco ponad górną linię
% koniec 

% debugowanie:
% %\iffalse %*** %\fi %***
% \tracingall przed problematycznym kodem - bardzo dużo informacji
% lualatex+se+debug - dodatkowe opcje: -file-line-error i -halt-on-error 
% \errorcontextlines
%\def\mydebug{1}
\begin{document}
%\makeatletter
%\grechangedim{interwordspacetext}{\unexpanded{\the\fontdimen2\font plus \the\gluestretch\gre@space@skip@interwordspacenotes minus \the\glueshrink\gre@space@skip@interwordspacenotes}}{fixed}%
%\makeatother
\rubryka{Odmawia się \emph{Pater noster} po cichu, oraz czyni rachunek sumienia. W przypadku obecności kapłana, może on poprowadzić rachunek sumienia, wówczas pierwszy odmawia  \emph{Confiteor}, a~następnie wszyscy (analogicznie jak przy modlitwach u stopni ołtarza).}%
\BeginParallel
\latN{\inicjal{P}{áter} nóster, qúi es in caélis, sanctificétur nómen túum: advéniat régnum túum: fíat volúntas túa, sícut in caélo et in térra. Pánem nóstrum quotidiánum da nóbis hódie: et dimítte nóbis débita nóstra, sícut et nos dimíttimus debitóribus nóstris: et ne nos indúcas in tentatiónem: sed líbera nos a málo. Amen.}%
\pol{\inicjal{O}{jcze} nasz, któryś jest w niebie, święć się imię Twoje: przyjdź królestwo Twoje: bądź wola Twoja, jako w niebie tak i na ziemi. Chleba naszego powszedniego daj nam dzisiaj: i odpuść nam nasze winy, jako i my odpuszczamy naszym winowajcom: i nie wódź nas na pokuszenie: ale nas zbaw ode złego. Amen.}%
\latN{\inicjal{C}{onfíteor} Déo omnipoténti, beátæ Maríæ semper Vírgini, beáto Michaéli Archángelo, beáto Joánni Baptístæ, sánctis Apóstolis Pétro et Páulo, et ómnibus Sánctis, quia peccávi nimis, cogitatióne, vérbo et ópere: méa cúlpa, méa cúlpa, méa máxima cúlpa. Ídeo précor beátam Maríam semper Vírginem, beátum Michaélem Archángelum, beátum Joánnem Baptístam, sánctos Apóstolos Pétrum et Páulum, et ómnes Sánctos, oráre pro me ad Dóminum Déum nóstrum.}%
\pol{\inicjal{S}{powiadam} się Bogu wszechmogącemu, Najświętszej Maryi zawsze Dziewicy, świętemu Michałowi Archaniołowi, świętemu Janowi Chrzcicielowi, świętym Apostołom Piotrowi i Pawłowi, i wszystkim Świętym, że bardzo zgrzeszyłem, myślą, mową i uczynkiem: bijąc się w piersi moja wina, moja wina, moja bardzo wielka wina. Przeto błagam Najświętszą Maryję zawsze Dziewicę, świętego Michała Archanioła, świętego Jana Chrzciciela, świętych Apostołów Piotra i Pawła, i wszystkich świętych, o modlitwę za mnie do Pana Boga naszego.}%
\latN{\inicjal{M}{isereátur} nóstri omnípotens Déus, et dimíssis peccátis nóstris, perdúcat nos ad vitam ætérnam. Amen.}%
\pol{\inicjal{N}{iech} się zmiłuje nad nami Bóg wszechmogący, a odpuściwszy nam grzechy, doprowadzi nas do życia wiecznego. Amen.}%
\latN{\inicjal{I}{ndulgéntiam}, \red{\grecross}~absolutiónem et remissiónem peccatórum nostrórum tríbuat nóbis omnípotens et miséricors Dóminus. Amen.}%
\pol{\inicjal{P}{rzebaczenia}, \red{\grecross}~rozgrzeszenia i odpuszczenia grzechów naszych, niech nam udzieli Pan wszechmogący i miłosierny. Amen.}%
\EndParallel
\end{document}